\subsection{\textit{Data Cleaning} dan \textit{Preprocessing}}
\label{subsec:nb-preprocessing}

\todo{Jelaskan \textit{pipeline} khusus NB beserta alasan tiap langkah, dikaitkan dengan temuan EDA (Subbab~\ref{sec:eda}) dan karakteristik algoritma. Semua transformasi di-\textit{fit} pada data latih saja lalu diterapkan ke data uji \cite{slide-dataprep}.}

\subsubsection{Penanganan \textit{Missing Values}}
\todo{NB dapat mengabaikan fitur yang hilang saat menghitung \textit{likelihood}, atau memakai imputasi. Jelaskan pilihan dan alasannya.}

\subsubsection{Data Duplikat dan Data Tidak Valid}
\todo{Jumlah duplikat, nilai tidak wajar, atau \textit{outlier} yang ditangani untuk NB dan caranya. Baris data uji tidak boleh dibuang.}

\subsubsection{Konstruksi Fitur dan \textit{Encoding}}
\todo{Fitur kategorikal dapat dipakai langsung (\textit{likelihood} frekuensi). Jelaskan apakah fitur numerik di-\textit{bin} (diskretisasi) atau dimodelkan Gaussian.}

\subsubsection{Normalisasi dan Standarisasi}
\todo{Tidak berpengaruh pada Gaussian NB (parameter diestimasi per fitur). Tulis alasannya singkat.}

\subsubsection{Penanganan \textit{Imbalanced Data}}
\todo{\textit{Prior} kelas sudah mencerminkan ketidakseimbangan. Jelaskan apakah \textit{prior} diubah atau \textit{threshold} disesuaikan.}

\subsubsection{\textit{Feature Selection}}
\todo{Penting karena fitur yang saling berkorelasi melanggar asumsi independensi dan menghitung bukti dua kali. Sebutkan fitur akhir yang dipakai.}

\subsubsection{\textit{Dimensionality Reduction}}
\todo{PCA menghasilkan komponen yang tidak berkorelasi sehingga lebih sesuai dengan asumsi independensi. Laporkan hasil percobaannya.}
